%% Section 2: the compiler and its languages. Budget: 1.5 pages
%% (notes/paper-outline.md, section 2). Grammars follow definitions.tex
%% (Types and Languages, Compilation) and agda/Mu/Language/*.agda; the Agda
%% module names match the paper's language names. Every pass but delay is an equation
%% (agda/Mu/Theorems/GraphModel.agda).

\section{The Compiler and Its Languages}
\label{sec:compiler}

\begin{figure}[t]
\centering
\small
\begin{tikzpicture}[x=3.6cm, y=1.0cm, every node/.style={inner sep=2pt},
    lbl/.style={font=\scriptsize}, ->, >=stealth]
  \node (s0) at (0,0)  {\iswim};
  \node (s1) at (1,0)  {\iswim};
  \node (c2) at (2,0)  {Clos};
  \node (c2b) at (3,0) {Clos};
  \node (c3) at (3,-1) {EnvClos};
  \node (c4) at (2,-1) {PairClos};
  \node (c5) at (1,-1) {GlobalFun};
  \node (mon) at (0,-1) {Mon};
  \node (imp) at (0,-2) {Imp};
  \node (cfg) at (1,-2) {CFG};
  \node (x86) at (2,-2) {x86$\mu$};
  \draw (s0) -- node[lbl,above] {\pass{simplify}} (s1);
  \draw (s1) -- node[lbl,above] {\pass{enclose}} (c2);
  \draw (c2) -- node[lbl,above] {\pass{optimize}} (c2b);
  \draw (c2b) -- node[lbl,right] {\pass{concretize}} (c3);
  \draw (c3) -- node[lbl,above] {\pass{delay}} (c4);
  \draw (c4) -- node[lbl,above] {\pass{globalize}} (c5);
  \draw (c5) -- node[lbl,above] {\pass{RCO}} (mon);
  \draw (mon) -- node[lbl,left] {\pass{lift locals}} (imp);
  \draw (imp) -- node[lbl,above] {\pass{explicate control}} (cfg);
  \draw (cfg) -- node[lbl,above] {\pass{select instructions}} (x86);
  \begin{scope}[on background layer]
    \draw[-, rounded corners=6pt, fill=gray!12, draw=gray!60, dashed]
      (1.25,0.42) rectangle (3.38,-1.3);
  \end{scope}
  \node[font=\itshape, text=gray!70!black] at (2.45,-0.5){closure conversion};
\end{tikzpicture}
\caption{The passes of the compiler and the languages between them.}
\label{fig:pipeline}
\end{figure}

\begin{figure}[t]
\small
\[
\begin{array}{@{}l@{\;\;}r@{\;}c@{\;}l@{}}
\text{Variables} & x, y, z & & \text{(from Imp on, slots of a stack frame)}\\
\text{Types} & A, C & ::= & B \mid A \rightarrow C \mid \kw{tup}(A_1, \ldots, A_n)
  \mid A + C \mid \mu X.\, A \mid X \mid \kw{code}(\Delta; A \rightarrow C)\\
 & \Delta & ::= & A_1, \ldots, A_n\\[3pt]
\iswim & K, L, M, N & ::= & x \mid \lambda x.\, N \mid L\;M \mid c
  \mid \kw{tuple}\;\vec{M} \mid \kw{get}_i\,M \mid \kw{inl}\;M \mid \kw{inr}\;M
  \mid \kw{case}\;L\;(x.\,M)\;(y.\,N)\\
 & & \mid & \kw{prim}_p\;L\;M \mid \kw{let}\;x = M\;\kw{in}\;N
  \mid \kw{fold}\;M \mid \kw{unfold}\;M\\[3pt]
\text{Clos} & K, L, M, N & ::= & \ldots \mid \kw{clos}\;(\vec{x}, y.\,N)\;\vec{M}
  \qquad (\text{in place of } \lambda x.\, N)\\
\text{EnvClos} & K, L, M, N & ::= & \ldots \mid \kw{clos}\;(z, y.\,N)\;\vec{M}\\
\text{PairClos} & K, L, M, N & ::= & \ldots \mid \kw{code}\;(z, y.\,N) \mid \kw{pair}\;L\;M \mid \kw{callc}\;L\;M
  \qquad (\text{in place of } \kw{clos} \text{ and } L\;M)\\
\text{GlobalFun} & K, L, M, N & ::= & \ldots \mid f_k
  \qquad (\text{in place of } \kw{code})\\
 & d & ::= & \kw{code}\;(z, y.\,N) \qquad P ::= (d_0, \ldots, d_{m-1};\, M)\\[3pt]
\text{Mon} & a & ::= & x \mid c \mid f_k\\
 & o & ::= & \kw{pair} \mid \kw{callc} \mid \kw{tuple} \mid \kw{get}_i \mid \kw{inl}
  \mid \kw{inr} \mid \kw{prim}_p \mid \kw{fold} \mid \kw{unfold}\\
 & M, N & ::= & a \mid o\;\vec{a} \mid \kw{case}\;a\;(x.\,M)\;(y.\,N)
  \mid \kw{let}\;x = M\;\kw{in}\;N
  \qquad P ::= (\vec{d};\, M)\\[3pt]
\text{Imp} & a & ::= & x \mid c \mid f_k
  \qquad r ::= a \mid o\;\vec{a}\\
 & e & ::= & r \mid x \mathrel{:=} e_1;\; e_2 \mid \kw{case}\;a\;(x.\,e_1)\;(x.\,e_2)
  \qquad P ::= (e_0, \ldots, e_{m-1};\, e)\\[3pt]
\text{CFG} & t & ::= & \kw{return}\;r \mid x \mathrel{:=} r;\; t \mid \kw{goto}\;l
  \mid \kw{case}\;a\;(x.\,l_1)\;(x.\,l_2)\\
 & b & ::= & (t_0, \ldots, t_{j-1};\, t)
  \qquad P ::= (\vec{b};\, b)\\[3pt]
\text{x86}\mu & a & ::= & x \mid c \mid f_k \mid \%r \mid \kw{fold}\;a \mid \kw{unfold}\;a
  \qquad d ::= \%r \mid x
  \qquad r ::= \kw{rax} \mid \kw{rdi} \mid \kw{rsi}\\
 & o & ::= & \kw{pair}(a_1, a_2) \mid \kw{tuple}(\vec{a}) \mid \kw{inl}(a) \mid \kw{inr}(a)\\
 & i & ::= & \kw{movq}\;a, d \mid \kw{subq}\;a, \%r \mid \kw{cmpq}\;a_2, a_1
  \mid \kw{sete}\;d \mid \kw{alloc}\;o, d \mid \kw{load}\;\kw{get}_i\;a, d\\
 & & \mid & \kw{fst}\;x, \%r \mid \kw{snd}\;x, \%r \mid \kw{callq}\;a\\
 & s & ::= & i;\; s \mid \kw{retq} \mid \kw{jmp}\;l \mid \kw{case}\;x\;(y.\,l_1)\;(y.\,l_2)\\
 & b & ::= & (n;\, s_0, \ldots, s_{j-1};\, s)
  \qquad P ::= (\vec{b};\, b)
\end{array}
\]
\caption{The grammars of the compiler's languages.}
\label{fig:grammars}
\end{figure}

Our compiler translates \iswim, a call-by-value $\lambda$-calculus
with recursive types, to x86$\mu$, a typed assembly language, in ten
passes through seven intermediate languages. Figure~\ref{fig:pipeline} shows the passes and languages, which this section defines and Section~\ref{sec:compiler-proof} proves correct.
Figure~\ref{fig:grammars} presents the grammars of all the languages, and Appendix~\ref{app:types} gives their typing rules.\footnote{Every language is intrinsically
typed in our Agda development, so every term is well typed by
construction and no pass needs a precondition about scoping or
typing.} In Figure~\ref{fig:grammars}, each closure language lists
only the forms it adds to the one above it, and ``$\ldots$'' stands
for the remaining forms of \iswim. Throughout, $c$ ranges over
constants, $p$ over the primitives $\kw{sub}$ and $\kw{eq}$, $l$ over
block labels, and $f_k$ over the names of a program's global
definitions, where $f_k$ refers to the definition $d_k$. In the
languages prior to Imp, $x$, $y$, and $z$ range over variables. From
Imp on, they range over fixed locations in a function's stack frame,
as we explain below. We write a binder as one or more names followed by a period, and the binder binds the names in the term after the period. For example, in
$\lambda x.\,N$ and $\kw{case}\;L\;(x.\,M)\;(y.\,N)$, the name $x$ is
bound in $N$ and $M$, and $\kw{clos}\;(\vec{x}, y.\,N)$ binds $\vec{x}$
and $y$ in $N$. Every language also has a graph-model semantics
(Section~\ref{sec:graph-models}), given in full in
Appendix~\ref{app:languages}.

\paragraph{The source language}

\iswim\ has integers, $n$-ary tuples, binary sums, the primitives
$\kw{sub}$ and $\kw{eq}$, and iso-recursive types. The primitive $\kw{eq}$ returns
$\kw{inl}\;\langle\rangle$ for true and $\kw{inr}\;\langle\rangle$ for
false, so $\kw{case}$ serves as the conditional. A recursive type $\mu
X.\, A$ is isomorphic to its unfolding $A[\mu X.\, A / X]$, and the
two coercions of this isomorphism are explicit. The $\kw{fold}$
operation turns a term of the unfolding into a term of the recursive
type, and $\kw{unfold}$ turns it back:
\[
  \frac{\Gamma \vdash M : A[\mu X.\, A / X]}{\Gamma \vdash \kw{fold}\;M : \mu X.\, A}
  \qquad
  \frac{\Gamma \vdash M : \mu X.\, A}{\Gamma \vdash \kw{unfold}\;M : A[\mu X.\, A / X]}
\]
Neither coercion does anything at run time. \iswim\ has no recursion
operator.  Recursive types make self-application typable, so a program can define the call-by-value fixed-point combinator:
\[
  \mathit{fix} = \lambda f.\, g\;(\kw{fold}\;g)
  \quad\text{where}\quad
  g = \lambda x.\, f\;(\lambda v.\, (\kw{unfold}\;x)\;x\;v)
\]
with $x : \mu X.\, X \rightarrow (\kw{Int} \rightarrow \kw{Int})$.

\paragraph{The optimizer}

The first pass, \pass{simplify} (\S\ref{app:pass-simplify}), is a simplified version of the
optimizer cp0 of Chez Scheme \citep{Waddell:1997fk}. It propagates
known values, folds constants, reduces a
$\kw{case}$, projection, or $\kw{unfold}$ of a known value, and
inlines small known functions.

\paragraph{Closure conversion}

The next five passes perform closure conversion. The \pass{enclose}
pass (\S\ref{app:pass-enclose}) turns each $\lambda$-abstraction into a closure
$\kw{clos}\;(\vec{x}, y.\,N)\;\vec{M}$ of Clos, which binds a name
$x_i$ for each captured value $M_i$ and the parameter $y$ in its code
$N$, so the code is closed. The \pass{enclose} pass captures every
variable in scope, and the \pass{optimize} pass (\S\ref{app:pass-optimize}) then drops the
captured values that the code does not use. The \pass{concretize} pass (\S\ref{app:pass-concretize})
stores the captured values of a closure in a single environment tuple,
so an EnvClos closure $\kw{clos}\;(z, y.\,N)\;\vec{M}$ binds two variables in its code, the environment $z$ and the parameter $y$, and
the code reads $x_i$ as $\kw{get}_i\,z$. The \pass{delay} pass (\S\ref{app:pass-delay}) then
separates a closure into its code and its environment, translating
from EnvClos to PairClos. In this target IR, $\kw{code}\;(z, y.\,N)$ is
closed code of type $\kw{code}(\Delta; A \rightarrow C)$, which
expects an environment of type $\kw{tup}(\Delta)$ and an argument of
type $A$. The form $\kw{pair}\;L\;M$ builds a closure, which is a pair of code
$L$ and an environment $M$ of the same $\Delta$. Its type, $A \rightarrow C$, hides $\Delta$, as in the existential type
$\exists \Delta.\, \kw{code}(\Delta; A \rightarrow C) \times
\kw{tup}(\Delta)$ of typed closure conversion \citep{Minamide:1996ys},
so closures with different environments have the same type. Thus
$\kw{pair}$ also packs an existential, but it means just a pair. The form
$\kw{callc}\;L\;M$ calls the code of the closure $L$ with its
environment and the argument $M$.
\vspace{-5pt}
\begin{gather*}
  \frac{z : \kw{tup}(\Delta),\, y : A \vdash N : C}
       {\Gamma \vdash \kw{code}\;(z, y.\,N) : \kw{code}(\Delta; A \rightarrow C)}
  \quad
  \frac{\begin{array}{l}
        \Gamma \vdash L : \kw{code}(\Delta; A \rightarrow C)\\
        \Gamma \vdash M : \kw{tup}(\Delta)
        \end{array}}
       {\Gamma \vdash \kw{pair}\;L\;M : A \rightarrow C}
  \quad
  \frac{\begin{array}{l}
        \Gamma \vdash L : A \rightarrow C \\
        \Gamma \vdash M : A
        \end{array}}
       {\Gamma \vdash \kw{callc}\;L\;M : C}
\end{gather*}
Finally, the \pass{globalize} pass (\S\ref{app:pass-globalize}) turns the code of every function into a
global definition $d_k$, and each occurrence of the code into a
reference $f_k$ to it, so a GlobalFun program is a list of function
definitions and a main term.

\paragraph{The back end}

The pass \pass{remove complex operands}
(\pass{RCO}, \S\ref{app:pass-rco}) binds each operand that is not atomic to a variable with
$\kw{let}$, so the operands of every operation in Mon are atomic.  The
atomic terms are variables, constants, and function names.  The
\pass{lift locals} pass (\S\ref{app:pass-lift-locals}) gives each variable a fixed location in the
function's stack frame, called a \emph{slot}, numbered by how many
variables are bound before it. The \pass{explicate control} pass (\S\ref{app:pass-explicate}) turns
the nested $\kw{case}$s of Imp into a control-flow graph. A CFG
program has a function body for each global definition and for the
main program. A body consists of numbered blocks and a starting tail,
and a block may jump only to blocks numbered before it, so control
flow is acyclic.

Finally, \pass{select instructions} (\S\ref{app:pass-select}) produces x86$\mu$, a typed assembly
language in the style of TAL \citep{Morrisett:1999fk} with the
registers $\kw{rax}$, $\kw{rdi}$, and $\kw{rsi}$, a frame of slots,
and an abstract heap. Programs have the same structure as in CFG, except that a function
body also declares the number of slots in its frame. An instruction reads operands $a$, which are slots, constants,
names of definitions, and registers, and writes to a destination $d$,
which is a register or a slot. The compiled code uses the registers to stage arithmetic and calls. The instruction $\kw{alloc}$ allocates an object
$o$: a closure ($\kw{pair}$), a tuple, or an injection. A call
$\kw{callc}\;x\;a$ of the closure in slot $x$ on the argument $a$
becomes the instruction sequence 
$\kw{fst}\;x, \%\kw{rax};\;\;
 \kw{snd}\;x, \%\kw{rdi};\;\;
 \kw{movq}\;a, \%\kw{rsi};\;\;
 \kw{callq}\;\%\kw{rax}
$.
The first two instructions load the code and the environment of the
closure, the third puts the argument in place, and the call leaves its
result in $\kw{rax}$. The type system
tracks which closure is open, so that a call never combines the code
of one closure with the environment of another.  The coercions $\kw{fold}$ and
$\kw{unfold}$ survive to x86$\mu$ as annotations on operands,
$\kw{fold}\;a$ and $\kw{unfold}\;a$, that do nothing.

x86$\mu$ abstracts from a real instruction set. It is typed; its heap is abstract, with no addresses or memory layout; it has
three registers and no register allocation; control flow within a
function is acyclic, so loops are recursive calls; and $\kw{case}$ is a
pseudo-instruction that stands for a test, a load, and a conditional
jump. None of these abstractions concerns closure
conversion, because in x86$\mu$, as in machine code, a closure is a heap object
that pairs the name of a function with a tuple of captured values, and
a call loads both before it jumps.
